%% =====================================================================
%%  comunicado.tex  --  MAESTRO del documento publico (plantilla)
%% ---------------------------------------------------------------------
%%  Un solo tipo para los cuatro generos publicos del despacho:
%%  COMUNICADO | PRONUNCIAMIENTO | NOTA DE PRENSA | ANUNCIO
%%  (la palabra de la cabecera sale de \setmeta{tipodoc} en metadata.tex).
%%
%%  Flujo de trabajo:
%%    1. Crear con numeracion automatica:
%%         scripts/nuevo.sh comunicado tema-breve --sumilla "Tema"
%%       (o pronunciamiento | nota-prensa | anuncio)
%%    2. Editar metadata.tex y el cuerpo de abajo.
%%    3. Compilar:  scripts/build.sh <carpeta>/comunicado.tex
%%
%%  Selecciones disponibles (descomentar para cambiar el default):
%%    \sinfotobanda               % banda sin retrato (temas sobrios)
%%    \sinconsignabanda           % banda sin el arte de la consigna
%%    \sinbandaprensa             % pie institucional normal (sin banda)
%%    \setbandafotoaltura{4.0cm}  % retrato mas discreto (defecto 4.50cm)
%%    \sinmarcaagua               % sin marca de agua
%%    \setinterlineado{1.05}      % compactar para caber en una hoja
%%    \setespacioparrafos{5pt}
%% =====================================================================
\documentclass{comunicado}

%% =====================================================================
%%  metadata.tex  --  DATOS del documento (lo unico que se edita aqui)
%% ---------------------------------------------------------------------
%%  El estilo lo pone el framework; este archivo solo declara datos.
%%  Los defaults del despacho (autor, cargo, lema, siglas, distrito)
%%  ya vienen del nucleo: sobrescribir solo lo que cambie.
%% =====================================================================

%% Genero del documento (cambia la palabra grande de la cabecera):
%%   Comunicado (defecto) | Pronunciamiento | Nota de Prensa | Anuncio
% \setmeta{tipodoc}{Pronunciamiento}

\setmeta{numero}{001}
\setmeta{anio}{2026}
% \setmeta{siglas}{ABC/CR}          % default del nucleo
%% Documento sin numero: comentar \setmeta{numero} (la linea desaparece)

\setmeta{lugar}{Lima}
\setmeta{fecha}{6 de agosto del 2026}   % comentar para usar la fecha de hoy

%% Nombre institucional del PDF (estandar del despacho):
%%    TIPO N°XXX-AAAA - Sumilla.pdf   (se compone al compilar)
%% Descomentar y poner la descripcion breve del tema:
% \setmeta{sumilla}{Encuentro por la educación ayacuchana}
%% O bien un nombre totalmente a medida:
% \setmeta{nombrearchivo}{COMUNICADO N°001-2026 - Encuentro educación}

%% Contacto de prensa propio (si difiere del despacho); lo usa
%% \contactoprensa en las notas de prensa:
% \setmeta{prensanombre}{Oficina de Prensa del Despacho}
% \setmeta{prensacorreo}{prensa@ejemplo.pe}
% \setmeta{prensacelular}{999 999 999}


\begin{document}

\cabeceradocumento

%% Titular y bajada: recomendados en notas de prensa; en comunicados
%% y pronunciamientos breves pueden omitirse.
\titular{El despacho convoca al Primer Encuentro por la Educación
  Ayacuchana}
\bajada{Autoridades, gremios docentes y comunidad educativa se reunirán
  para articular una agenda regional por la escuela pública.}

El despacho del \meta{cargo} \meta{autor} se dirige a la opinión
pública, a la comunidad educativa ayacuchana y a los medios de
comunicación para expresar lo siguiente:

\begin{puntos}
  \item La educación pública de Ayacucho atraviesa una situación que
    exige respuestas articuladas entre el Estado, los gremios y la
    sociedad civil: brechas de infraestructura, deuda social con los
    maestros y desigualdades que golpean con mayor dureza a la
    escuela rural.

  \item En ejercicio de la función de representación, este despacho
    convoca al \textbf{Primer Encuentro por la Educación Ayacuchana},
    que se realizará el sábado 8 de agosto del 2026, a las 9:00 a.\,m.,
    en el auditorio de la Casa del Maestro (Jr.~2 de Mayo N.º~235,
    Ayacucho).

  \item Los acuerdos del encuentro serán recogidos en una agenda
    legislativa regional que este despacho impulsará en la Cámara de
    Diputados, con seguimiento público de sus avances.
\end{puntos}

Reafirmamos nuestro compromiso con una educación pública, gratuita y de
calidad al servicio del pueblo ayacuchano.

\cierrelugarfecha

%% Los documentos publicos van SIN firma ni posfirma (pedido de Edison,
%% 2026-08-06): la voz institucional —nombre, cargo y foto del diputado—
%% ya vive en la banda de prensa al pie. Si un pronunciamiento formal
%% necesitara cierre firmado, existe \firmainstitucional (structures/).
%%
%% Solo NOTAS DE PRENSA: bloque de contacto para los medios.
% \contactoprensa

\end{document}
